% Part: computability
% Chapter: recursive-functions
% Section: sequences

\documentclass[../../../include/open-logic-section]{subfiles}

\begin{document}

\olfileid{cmp}{rec}{seq}
\olsection{سلسلے}

بنیادی بازگشتی فنکشناں دا سیٹ خاصا مضبوط اے۔ پر جدوں ساڈے کول
\emph{سلسلیاں} نال کم کرن دا مناسب طریقہ ہووے گا، تاں اسی ہور
وی زیادہ کم کر سکاں گے۔ اسی قدرتی عدداں دے محدود سلسلیاں
نوں قدرتی عدداں نال ایویں جوڑاں گے: سلسلہ
$\langle a_0, a_1, a_2, \dots, a_k \rangle$ ایس عدد دے
مطابق اے:
\[
p_0^{a_0+1} \cdot p_1^{a_1+1} \cdot p_2^{a_2+1} \cdot \dots \cdot
p_k^{a_k+1}.
\]
اسی ہر اسّ وچ اک شامل کردے آں تاکہ، مثال لئی، سلسلے
$\langle 2, 7, 3\rangle$ تے $\langle 2, 7, 3, 0, 0 \rangle$
دے عددی کوڈ وکھرے ہون۔ $0$ تے~$1$ دونوں نوں خالی
سلسلے دا کوڈ منیا جا سکدا اے؛ پکی چُون لئی
$\emptyseq$ توں مراد~$0$ لاں گے۔

سلسلیاں دی ایہہ کوڈکاری اس کرکے کم کردی اے کہ حساب دا
بنیادی قضیہ کہندا اے: ہر قدرتی عدد $n \ge 2$ نوں اکّو
اک طریقے نال ایس روپ وچ لکھیا جا سکدا اے:
\[
n = p_0^{a_0} \cdot p_1^{a_1} \cdot \dots \cdot p_k^{a_k}
\]
جتھے $a_k \ge 1$ ہووے۔ ایس توں پکا ہوندا اے کہ نقشہ
$\tuple{}(a_0,
\dots, a_k) = \tuple{a_0, \dots, a_k}$ اک بہ اک اے:
وکھرے سلسلے وکھرے عدداں نال جُڑدے نیں؛ ہر عدد نال
زیادہ توں زیادہ اک سلسلہ جُڑدا اے۔

ہُن اسی وکھاواں گے کہ سلسلے دی لمبائی لبھنا، اوہدا
$i$واں رکن لبھنا، اک رکن سلسلے دے آخر شامل کرنا، تے
دو سلسلے جوڑنا، سب بنیادی بازگشتی عمل نیں۔

\begin{prop}
فنکشن $\len{s}$، جیہڑا سلسلے $s$ دی لمبائی واپس
دیندا اے، بنیادی بازگشتی اے۔
\end{prop}

\begin{proof}
فرض کرو $R(i, s)$ ایہہ تعلق اے:
\[
R(i, s) \text{ اوہدوں تے صرف اوہدوں جدوں }
p_i \mid s \land p_{i+1} \nmid s.
\]
$R$ واضح طور اُتے بنیادی بازگشتی اے۔ جدوں وی $s$
غیر خالی سلسلے دا کوڈ ہووے، یعنی
\[
s = p_0^{a_0+1} \cdot \dots \cdot p_{k}^{a_{k}+1},
\]
تاں $R(i,s)$ اوہدوں درست اے جدوں $p_i$ سب توں
وڈا اولی عدد ہووے جیہدے لئی $p_i \mid s$؛
یعنی $i = k$۔ ایس لئی
$s$ دی لمبائی $i+1$ اوہدوں تے صرف اوہدوں اے
جدوں $p_i$ سب توں وڈا اولی عدد ہووے جیہڑا~$s$
نوں پورا ونڈدا اے۔ ایس کرکے اسی ایویں رکھ سکدے آں:
\[
\len{s} =
\begin{cases}
  0 & \text{جے $s = 0$ یا $s = 1$ ہووے} \\
  1 + \bmin{i < s}{R(i, s)} & \text{ورنہ}
\end{cases}
\]
ایتھے محدود کمترین قدر دی تلاش ورت سکدے آں؛ کیونکہ
صرف اک $i$ ای $R(i,s)$ نوں پورا کردا اے، جے
$s$~سلسلے دا کوڈ ہووے؛ تے جے $i$ موجود
ہووے تاں اوہ خود~$s$ توں چھوٹا ہوندا اے۔
\end{proof}

\begin{prop}
فنکشن $\fn{append}(s,a)$، جیہڑا $a$ نوں سلسلے
$s$ دے آخر شامل کرن دا نتیجہ واپس دیندا اے، بنیادی بازگشتی اے۔
\end{prop}

\begin{proof}
$\fn{append}$ نوں ایویں متعین کیتا جا سکدا اے:
\[
\fn{append}(s,a) =
\begin{cases}
  2^{a+1} & \text{جے $s = 0$ یا $s = 1$ ہووے} \\
  s \cdot p_{\len{s}}^{a+1} & \text{ورنہ}
\end{cases}
\]
\end{proof}

\begin{prop}
فنکشن $\fn{element}(s,i)$، جیہڑا $i$واں رکن
سلسلے $s$ وچوں
واپس دیندا اے (شروع دے رکن نوں $0$واں آکھدے آں)،
یا $0$ واپس دیندا اے جے $i$، $s$ دی لمبائی دے
برابر یا اوہدے توں وڈا ہووے، بنیادی بازگشتی اے۔
\end{prop}

\begin{proof}
غور کرو کہ $a$، $i$واں رکن~$s$ دا اوہدوں تے
صرف اوہدوں اے جدوں $p_i^{a+1}$،~$p_i$ دی اوہ
سب توں وڈی قوت ہووے جیہڑی~$s$ نوں پورا ونڈدی
اے؛ یعنی $p_i^{a+1} \mid s$ پر
$p_i^{a+2} \nmid s$۔ ایس لئی:
\[
\fn{element}(s,i) =
\begin{cases}
  0 & \mbox{جے $i \geq \len{s}$ ہووے} \\
  \bmin{a < s}{(p_i^{a+2} \nmid s)} & \text{ورنہ}
\end{cases}
\]
\end{proof}

اوپر متعین فنکشناں دے باقاعدہ ناں ہر وار ورتن دی بجائے
اسی مختصر علامتاں لیانیاں آں۔ $(s)_i$ نوں
$\fn{element}(s,i)$ دی جگہ لکھاں گے، تے $\tuple{s_0, \dots, s_k}$
نال ایہہ مختصر کراں گے:
\[
\fn{append}(\fn{append}(\dots \fn{append}(\emptyseq,s_0)
\dots),s_k).
\]
غور کرو کہ جے $s$ دی لمبائی~$k$ اے، تاں
$s$ دے رکن $(s)_0$، \dots،~$(s)_{k-1}$ نیں۔

\begin{prop}
فنکشن $\fn{concat}(s,t)$، جیہڑا دو سلسلے
جوڑدا اے، بنیادی بازگشتی اے۔
\end{prop}

\begin{proof}
سانوں ایہو جیہا فنکشن $\fn{concat}$ چاہیدا اے جیہدے لئی
\[
  \fn{concat}(\tuple{a_0, \dots, a_k}, \tuple{b_0, \dots, b_l}) =
  \tuple{a_0, \dots, a_k, b_0, \dots, b_l}.
\]
اسی اک ``مددگار'' فنکشن
$\fn{hconcat}(s,t,n)$ ورتاں گے جیہڑا پہلے $n$ رکن
$t$ توں لے کے~$s$ دے آخر جوڑدا اے۔ ایہہ فنکشن بنیادی
بازگشت راہیں ایویں متعین ہو سکدا اے:
\begin{align*}
  \fn{hconcat}(s,t,0) & = s\\
  \fn{hconcat}(s,t,n+1) & = \fn{append}(\fn{hconcat}(s,t,n),(t)_n)
  \intertext{فیر اسی $\fn{concat}$ نوں ایویں متعین کر سکدے آں:}
  \fn{concat}(s,t) & = \fn{hconcat}(s,t,\len{t}).
\end{align*}
\end{proof}

اسی $s \concat t$ نوں $\fn{concat}(s,t)$ دی بجائے لکھاں گے۔

ساڈے لئی سلسلے دے عددی کوڈ دی اک حد، اوہدی لمبائی
تے سب توں وڈے رکن دے حساب نال لبھ سکنا مفید ہووے گا۔
فرض کرو $s$ اک سلسلہ اے جیہدی لمبائی~$k$ اے، تے
ہر رکن کسے عدد~$x$ توں چھوٹا یا اوہدے برابر اے۔
تاں $s$ دے زیادہ توں زیادہ $k$ اولی عوامل نیں،
ہر اک زیادہ توں زیادہ~$p_{k-1}$ اے، تے ہر اک دا
اسّ زیادہ توں زیادہ $x+1$ ہوندا اے، جدوں~$s$ دی اولی عامل
بندی لکھی جاندی اے۔ دوجے لفظاں وچ، مثبت لمبائی لئی اسی
ایویں متعین کردے آں؛ خالی سلسلے لئی حد اک مقرر اے:
\[
\fn{sequenceBound}(x,k) = p_{k-1}^{k \cdot (x+1)},
\]
تاں اوپر دسے سلسلے~$s$ دا عددی کوڈ زیادہ توں
زیادہ~$\fn{sequenceBound}(x,k)$ اے۔

سلسلیاں دی ایہہ حد سانوں محدود تلاش نال نویں فنکشن
متعین کرن دا طریقہ دیندی اے۔ مثال لئی، اسی
$\fn{concat}$ نوں محدود تلاش نال متعین کر سکدے آں۔
سانوں صرف لبھّن والی شے (جوڑے گئے سلسلے دے عدد)
دی اک بنیادی بازگشتی \emph{تعیین} تے ایہہ حد دینی
اے کہ کتھے تک لبھنا اے۔ ایہہ تعریف کم کردی اے:
\begin{align*}
  \fn{concat}(s,t) = {} & \bmin{v < \fn{sequenceBound}(s+t,\len{s} +
    \len{t})}{} \\
  & \quad(\len{v} = \len{s} + \len{t} \land {}\\
    & \qquad \bforall{i < \len{s}}{((v)_i = (s)_i)} \land {} \\
      & \qquad \bforall{j < \len{t}}{((v)_{\len{s}+j} = (t)_j)})
\end{align*}

\begin{prob}
وکھاؤ کہ اک بنیادی بازگشتی فنکشن~$\fn{sconcat}(s)$
موجود اے جیہدے لئی
\[
\fn{sconcat}(\tuple{s_0, \dots, s_k}) = s_0 \concat \dots \concat s_k.
\]
\end{prob}

\begin{prob}
وکھاؤ کہ اک بنیادی بازگشتی فنکشن~$\fn{tail}(s)$
موجود اے جیہدے لئی
\begin{align*}
  \fn{tail}(\emptyseq) & = 0 \text{ تے}\\
  \fn{tail}(\tuple{s_0, \dots, s_{k}}) & = \tuple{s_1, \dots, s_{k}}.
\end{align*}
\end{prob}

\begin{prop}
  \ollabel{prop:subseq}
فنکشن $\fn{subseq}(s, i, n)$، جیہڑا $s$ دا اوہ
ذیلی سلسلہ واپس دیندا اے جیہدی لمبائی~$n$ اے تے
جیہڑا $i$ویں رکن توں شروع ہوندا اے، بنیادی بازگشتی اے۔
\end{prop}

\begin{proof}
مشق۔
\end{proof}

\begin{prob}
  \olref[cmp][rec][seq]{prop:subseq} ثابت کرو۔
\end{prob}
  
\end{document}
